% GENERATED FILE. Edit canonical lessons, then run npm run generate:books.
% canonical-source-hash: fnv1a64:83ad3e2639d5aada
% canonical-lessons: SA-C331-bhaktih, SA-C331-candanam, SA-C331-visarjayati, SA-C331-japati, SA-C331-pavitrah

\chapter{Down to the River}
\label{ch:sa-down-to-the-river}
\hlchaptermodality{331}

\begin{chapteropening}
\textbf{By the end of this chapter:} \emph{I can say what people did at the river on a festival.}

Complete the last lesson of chapter 331: I can say what people did at the river on a festival.
\end{chapteropening}

\section[\texorpdfstring{\sk{भक्तिः} (bhaktiḥ)}{भक्तिः (bhaktiḥ)}]{\sk{भक्तिः} (bhaktiḥ) — devotion}

\label{lesson:SA-C331-bhaktih}



\begin{quote}
Before the new one: say the Sanskrit for a festival day, then the Sanskrit for a forehead mark.
\end{quote}

\subsection*{You'll want to know: \sk{भक्तिः}}
\textbf{\sk{भक्तिः}} — \emph{bhaktiḥ} — \textquotedblleft{}devotion\textquotedblright{}.

\begin{quote}
\sk{जनाः} \sk{भक्त्या} \sk{नदीं} \sk{नमन्ति।}
\end{quote}

\begin{quote}
\emph{janāḥ bhaktyā nadīṃ namanti.}
\end{quote}

\begin{quote}
People bow to the river with devotion.
\end{quote}

A feminine noun.

\begin{cousinweb}
\textbf{\sk{भक्तिः}} (\emph{bhaktiḥ}) comes from the root √bhaj, \textquotedblleft{}to share, worship\textquotedblright{}, as in \emph{bhajati}; MW gives \emph{bhakti} as a share, and also attachment, devotion, worship, faith (MW).
\end{cousinweb}

\begin{grammarlens}[title={What you've built}]
One more word for telling what happened.
\end{grammarlens}

\subsection*{Your turn}
\begin{itemize}
\raggedright
  \item \emph{Say it:} \emph{bhaktiḥ}
  \item \emph{Say it:} \emph{bhaktiḥ}, once more
  \item \emph{Say it:} \emph{janāḥ bhaktyā nadīṃ namanti.} — People bow to the river with devotion.
  \item \emph{Recall:} say the Sanskrit for a festival day, then the Sanskrit for a forehead mark, then say \emph{bhaktiḥ} again
\end{itemize}

\begin{tcolorbox}[breakable,colback=teal!4,colframe=teal!35!black,title={Before you move on}]
What does \sk{भक्तिः} mean? (\textquotedblleft{}Devotion\textquotedblright{}.) Say it once more.
\end{tcolorbox}
\section[\texorpdfstring{\sk{चन्दनम्} (candanam)}{चन्दनम् (candanam)}]{\sk{चन्दनम्} (candanam) — sandalwood; sandal paste}

\label{lesson:SA-C331-candanam}



\begin{quote}
Before the new one: say the Sanskrit for a forehead mark, then the Sanskrit for devotion.
\end{quote}

\subsection*{You'll want to know: \sk{चन्दनम्}}
\textbf{\sk{चन्दनम्}} — \emph{candanam} — \textquotedblleft{}sandalwood; sandal paste\textquotedblright{}.

\begin{quote}
\sk{जनाः} \sk{ललाटे} \sk{चन्दनं} \sk{धारयन्ति।}
\end{quote}

\begin{quote}
\emph{janāḥ lalāṭe candanaṃ dhārayanti.}
\end{quote}

\begin{quote}
People wear sandal paste on the forehead.
\end{quote}

A neuter noun.

\begin{cousinweb}
\textbf{\sk{चन्दनम्}} (\emph{candanam}) MW gives \emph{candana} as sandal, the tree, the wood or the fragrant paste made from it, without a derivation, so its origin is uncertain (MW).
\end{cousinweb}

\begin{grammarlens}[title={What you've built}]
One more word for telling what happened.
\end{grammarlens}

\subsection*{Your turn}
\begin{itemize}
\raggedright
  \item \emph{Say it:} \emph{candanam}
  \item \emph{Say it:} \emph{candanam}, once more
  \item \emph{Say it:} \emph{janāḥ lalāṭe candanaṃ dhārayanti.} — People wear sandal paste on the forehead.
  \item \emph{Recall:} say the Sanskrit for a forehead mark, then the Sanskrit for devotion, then say \emph{candanam} again
\end{itemize}

\begin{tcolorbox}[breakable,colback=teal!4,colframe=teal!35!black,title={Before you move on}]
What does \sk{चन्दनम्} mean? (\textquotedblleft{}Sandalwood; sandal paste\textquotedblright{}.) Say it once more.
\end{tcolorbox}
\section[\texorpdfstring{\sk{विसर्जयति} (visarjayati)}{विसर्जयति (visarjayati)}]{\sk{विसर्जयति} (visarjayati) — immerses; lets go}

\label{lesson:SA-C331-visarjayati}



\begin{quote}
Before the new one: say the Sanskrit for devotion, then the Sanskrit for sandalwood.
\end{quote}

\subsection*{You'll want to know: \sk{विसर्जयति}}
\textbf{\sk{विसर्जयति}} — \emph{visarjayati} — \textquotedblleft{}immerses; lets go\textquotedblright{}.

\begin{quote}
\sk{उत्सवस्य} \sk{अन्ते} \sk{जनाः} \sk{दीपान्} \sk{नद्यां} \sk{विसर्जयन्ति।}
\end{quote}

\begin{quote}
\emph{utsavasya ante janāḥ dīpān nadyāṃ visarjayanti.}
\end{quote}

\begin{quote}
At the end of the festival people set lamps afloat on the river.
\end{quote}

Past (he …): \sk{सः} \sk{दीपं} \sk{नद्यां} \sk{विसर्जितवान्}, \emph{saḥ dīpaṃ nadyāṃ visarjitavān}.

\begin{cousinweb}
\textbf{\sk{विसर्जयति}} (\emph{visarjayati}) is built from \emph{vi} and the causative of the root √sṛj, \textquotedblleft{}to let go, send forth\textquotedblright{}; MW gives \emph{vi-sṛj} as to send or pour forth, let go, and the immersion of an image is its festival use (MW).
\end{cousinweb}

\begin{grammarlens}[title={What you've built}]
One more word for telling what happened.
\end{grammarlens}

\subsection*{Your turn}
\begin{itemize}
\raggedright
  \item \emph{Say it:} \emph{visarjayati}
  \item \emph{Say it:} \emph{visarjayati}, once more
  \item \emph{Say it:} \emph{utsavasya ante janāḥ dīpān nadyāṃ visarjayanti.} — At the end of the festival people set lamps afloat on the river.
  \item \emph{Recall:} say the Sanskrit for devotion, then the Sanskrit for sandalwood, then say \emph{visarjayati} again
\end{itemize}

\begin{tcolorbox}[breakable,colback=teal!4,colframe=teal!35!black,title={Before you move on}]
What does \sk{विसर्जयति} mean? (\textquotedblleft{}Immerses; lets go\textquotedblright{}.) Say it once more.
\end{tcolorbox}
\section[\texorpdfstring{\sk{जपति} (japati)}{जपति (japati)}]{\sk{जपति} (japati) — chants softly; recites a prayer}

\label{lesson:SA-C331-japati}



\begin{quote}
Before the new one: say the Sanskrit for sandalwood, then the Sanskrit for immerses.
\end{quote}

\subsection*{You'll want to know: \sk{जपति}}
\textbf{\sk{जपति}} — \emph{japati} — \textquotedblleft{}chants softly; recites a prayer\textquotedblright{}.

\begin{quote}
\sk{वृद्धः} \sk{नदीतीरे} \sk{मन्त्रं} \sk{जपति।}
\end{quote}

\begin{quote}
\emph{vṛddhaḥ nadītīre mantraṃ japati.}
\end{quote}

\begin{quote}
The old man softly chants a mantra on the riverbank.
\end{quote}

Past (he …): \sk{सः} \sk{मन्त्रं} \sk{जपितवान्}, \emph{saḥ mantraṃ japitavān}.

\begin{cousinweb}
\textbf{\sk{जपति}} (\emph{japati}) comes from the root √jap, \textquotedblleft{}to utter in a low voice, mutter prayers\textquotedblright{} (MW).
\end{cousinweb}

\begin{grammarlens}[title={What you've built}]
One more word for telling what happened.
\end{grammarlens}

\subsection*{Your turn}
\begin{itemize}
\raggedright
  \item \emph{Say it:} \emph{japati}
  \item \emph{Say it:} \emph{japati}, once more
  \item \emph{Say it:} \emph{vṛddhaḥ nadītīre mantraṃ japati.} — The old man softly chants a mantra on the riverbank.
  \item \emph{Recall:} say the Sanskrit for sandalwood, then the Sanskrit for immerses, then say \emph{japati} again
\end{itemize}

\begin{tcolorbox}[breakable,colback=teal!4,colframe=teal!35!black,title={Before you move on}]
What does \sk{जपति} mean? (\textquotedblleft{}Chants softly; recites a prayer\textquotedblright{}.) Say it once more.
\end{tcolorbox}
\section[\texorpdfstring{\sk{पवित्रः} (pavitraḥ)}{पवित्रः (pavitraḥ)}]{\sk{पवित्रः} (pavitraḥ) — holy; pure}

\label{lesson:SA-C331-pavitrah}



\begin{quote}
Before the new one: say the Sanskrit for immerses, then the Sanskrit for chants softly.
\end{quote}

\subsection*{You'll want to know: \sk{पवित्रः}}
\textbf{\sk{पवित्रः}} — \emph{pavitraḥ} — \textquotedblleft{}holy; pure\textquotedblright{}.

\begin{quote}
\sk{सा} \sk{नदी} \sk{पवित्रा} \sk{अस्ति।}
\end{quote}

\begin{quote}
\emph{sā nadī pavitrā asti.}
\end{quote}

\begin{quote}
That river is holy.
\end{quote}

\begin{cousinweb}
\textbf{\sk{पवित्रः}} (\emph{pavitraḥ}) MW gives \emph{pavitra} first as a means of purification, a filter, and as an adjective purifying, pure, holy, sacred, without a root in the entry, so its origin is left uncertain here (MW).
\end{cousinweb}

\begin{grammarlens}[title={What you've built}]
One more word for telling what happened.
\end{grammarlens}

\subsection*{Your turn}
\begin{itemize}
\raggedright
  \item \emph{Say it:} \emph{pavitraḥ}
  \item \emph{Say it:} \emph{pavitraḥ}, once more
  \item \emph{Say it:} \emph{sā nadī pavitrā asti.} — That river is holy.
  \item \emph{Recall:} say the Sanskrit for immerses, then the Sanskrit for chants softly, then say \emph{pavitraḥ} again
\end{itemize}

\begin{tcolorbox}[breakable,colback=teal!4,colframe=teal!35!black,title={Before you move on}]
What does \sk{पवित्रः} mean? (\textquotedblleft{}Holy; pure\textquotedblright{}.) Say it once more.
\end{tcolorbox}
